

\documentclass{article}
\usepackage[utf8]{inputenc}
\usepackage{authblk}
\usepackage{setspace}
\usepackage[margin=1.25in]{geometry}
\usepackage{graphicx}
\graphicspath{ {./figures/} }
\usepackage{subcaption}
\usepackage{amsmath}
\usepackage{lineno}
\linenumbers

%%%%%% Bibliography %%%%%%
% Replace "sample" in the \addbibresource line below with the name of your .bib file.
\usepackage[style=ieee, 
citestyle=numeric-comp,
sorting=none]{biblatex}
\addbibresource{sample.bib}

%%%%%% Title %%%%%%
% Full titles can be a maximum of 200 characters, including spaces. 
% Title Format: Use title case, capitalizing the first letter of each word, except for certain small words, such as articles and short prepositions
\title{BME Frontiers Research Article \LaTeX\ Template}

%%%%%% Authors %%%%%%
% Authors should be listed in order of contribution to the paper, by first name, then middle initial (if any), followed by last name.
% Authors should be listed in the order in which they will appear in the published version if the manuscript is accepted. 
% Use an asterisk (*) to identify the corresponding author, and be sure to include that person’s e-mail address. Use symbols (in this order: †, ‡, §, ||, ¶, #, ††, ‡‡, etc.) for author notes, such as present addresses, “These authors contributed equally to this work” notations, and similar information.
% You can include group authors, but please include a list of the actual authors (the group members) in the Supplementary Materials.
\author[1*$\dag$]{Author One}
\author[2$\dag$]{Author Two}
\author[2]{Author Three}
\author[1,2]{Author Four}

%%%%%% Affiliations %%%%%%
\affil[1]{Department of Physics, A University, City, Country.}
\affil[2]{Department of Astronomy, B University, City, Country.}
\affil[*]{Corresponding author. Email: email@email.com}
\affil[$\dag$]{These authors contributed equally to this work.}

%%%%%% Date %%%%%%
% Date is optional
\date{}

%%%%%% Spacing %%%%%%
% Use paragraph spacing of 1.5 or 2 (for double spacing, use command \doublespacing)
\onehalfspacing

\begin{document}

\maketitle

%%%%%% Abstract %%%%%%
\begin{abstract}
The abstract should be a single paragraph written in plain language that a general reader can understand. Use a Structured Abstract with the sections indicated below. Do not include citations or undefined abbreviations in the abstract. Any abbreviations that appear in the title should be defined in the abstract. The length should not exceed 250 words. The sections to include are: 
\begin{itemize}
    \item Objective: An opening sentence that states the objective of the research
    \item Impact Statement: Brief description about the novelty and impact of the research
    \item Introduction: Enough background information to give context to the study
    \item Methods: A brief statement of the primary methods used by the study
    \item Results: A brief statement of primary results
    \item Conclusion: A short concluding sentence of the main take-home point(s) of the study
\end{itemize} 
\end{abstract}

%%%%%% Main Text %%%%%%

\section{Introduction}
Your manuscript should contain all of the numbered sections specified in this template: Introduction, Results, Discussion, Materials and Methods.

The manuscript should start with a brief introduction that lays out the problem addressed by the research and describes the paper’s importance. The scientific question being investigated should be described in detail. The introduction should provide sufficient background information to make the article understandable to readers in other disciplines and provide enough context to ensure that the implications of the experimental findings are clear. 

%%%%% Citations in the text %%%%%%
\subsection*{Citations}
Citations of references in the text should be identified using numbers in square brackets e.g., ``as discussed by Cui \cite{Cui1}'' or ``as discussed elsewhere \cite{Cui1,Ninomiya1,Li1,Wang1,Yang1}.'' All references should be cited within the text and uncited references will be removed. 

As an example, this template includes a ``sample.bib'' file containing the references in BibTeX.

%%%%%% Equations %%%%%%
\subsection*{Equations}
Equations should be provided in a text format, rather than as an image. Equations should be numbered consecutively, in round brackets, on the right-hand side of the page by using the ``\textbackslash begin\{equation\}'' command. They should be referred to as Equation 1, etc. in the main text.

\medskip For example, see Equation \ref{eq:1} and Equation \ref{eq:2} below.
\begin{equation} \label{eq:1}
    a^2 + b^2 = c^2
\end{equation}
\begin{equation} \label{eq:2}
\begin{split}
A & = \frac{\pi r^2}{2} \\
 & = \frac{1}{2} \pi r^2
\end{split}
\end{equation}

%%%%%% Figures %%%%%%
\subsection*{Figures}
Figures should be called out within the text and numbered in the order of their citation in the text. Every figure must have a descriptive title beginning with ``Figure [Number] …'' All figure titles should be either a phrase or a sentence; do not mix the two styles. See Figure \ref{fig:1} for example.
\begin{figure}[h]
    \centering
    \includegraphics[width=0.5\textwidth]{fig 1}
    \caption{This is an example figure.}
    \label{fig:1}
\end{figure}

Figures should be displayed on a white background. When preparing figures, consider that they can occupy either a single column (half page width) or two columns (full page width), and should be sized accordingly.

If a figure consists of multiple panels, they should be ordered logically and labelled with lower case roman letters (i.e., a, b, c, etc.). All labels should be explained in the legend. See Figure \ref{fig:2} for example.

Upon acceptance, authors will be asked to provide the figures as separate electronic files. At that stage, figures should be supplied in either vector art formats (PS, EPS, FIG, AI, Visio, WMF, EMF, Word, Excel, PowerPoint, OPJ, CDR, or PDF) or bitmap formats (Photoshop, TIFF, GIF, JPEG, PNG, BMP, etc.). Bitmap (BMP) images should be of at least 300 dpi resolution, unless due to the limited resolution of a scientific instrument. If a bitmap image has labels, the image and labels should be embedded in separate layers.
\begin{figure}[h]
    \centering
    \begin{subfigure}{0.4\textwidth}
        \includegraphics[width=0.9\textwidth, height=2in]{fig 1}
        \caption{\label{fig:2a}}
    \end{subfigure}
    \begin{subfigure}{0.4\textwidth}
        \includegraphics[width=0.9\textwidth, height=2in]{fig 2}
        \caption{\label{fig:2b}}
    \end{subfigure}
    \caption{This is an example of a figure consisting of multiple panels.     (\subref{fig:2a}) This is the first panel. (\subref{fig:2b}) This is the second panel.}
    \label{fig:2}
\end{figure}

%%%%%% Tables %%%%%%
\subsection*{Tables}
Tables should supplement, not duplicate, the text. They should be called out consecutively within the text and numbered in the order of their citation in the text. 

Every table must have a descriptive title beginning with ``Table [Number] …'' as noted in Table \ref{tab:1}. If numerical measurements are given, the units should be included in the column heading. Every vertical column should have a heading, followed by a unit of measure (if any) in parentheses. Units should not change within a column. Vertical rules should not be used. 

Centered headings of the body of the table can be used to break the entries into groups. Do not use footnotes in column heads; include any such details in sentence form in the table legend. Footnotes should contain information relevant to specific cells of the table; use the following symbols in order, as needed: $*, \dag, \ddag, \S, \|, \P, \#, **, \dag\dag$, etc.

\begin{table}[b]
    \caption{This is an example table.}    
    \centering
    \begin{tabular}{ccc}
            \hline
            Column 1 & Column 2 & Column 3 \\  
            \hline
            Cell 1 & Cell 2 & Cell 3\\ 
            Cell 4 & Cell 5 & Cell 6 \\
            \hline
            \end{tabular}

    \label{tab:1}
\end{table}

\section{Results}
The results should describe the experiments performed and the findings observed. The results section should be divided into subsections to delineate different experimental themes. Subheadings should either be all phrases or all complete sentences. All data must be shown either in the main text or in the Supplementary Materials.  
\begin{itemize}
    \item All data, both that in the main part of the manuscript and that in the Supplementary Materials, should be presented in the Results. No data should be presented for the first time in the Discussion. Data (such as from Western blots) should be appropriately quantified from multiple independent experiments. Inclusion of biological data from individual experiments that have not been repeated at least twice is generally not permitted.
    \item Subheadings must be either all complete sentences or all phrases. They should be brief, ideally less than 10 words. Subheadings should not end in a period. The Results section may have as many level 1 subheadings as are necessary.
    \item Figures and tables must be called out in numerical order. For example, the first mention of any panel of Figure 3 cannot precede the first mention of all panels of Figure 2. The supplementary figures (for example, Figure S1) and tables (Table S1) must also be called out in numerical order. All figures and tables should have a title and a legend.
    \item Mathematical expressions within a sentence of text should be created with ordinary Word characters; if this is not possible, then use MathType (or the equivalent). Only use MathType when necessary — for example, characters with overbars or carets, with stacked superscripts and subscripts, or within square root symbols.
\end{itemize}

\section{Discussion}
Include a Discussion that summarizes (but does not merely repeat) your conclusions and elaborates on their implications. There should be a paragraph outlining the limitations of your results and interpretation, as well as a discussion of the steps that need to be taken for the findings to be applied. Please avoid claims of priority. 

\section{Materials and Methods}
The materials and methods section should provide sufficient information to allow replication of the results. This section should be broken up by subheadings. Under exceptional circumstances, when a particularly lengthy description is required, a portion of the materials and methods can be included in the Supplementary Materials. 

\subsection{Experimental and Technical Design}
Begin the Materials and Methods with a subsection titled “Experimental and Technical Design” describing the objectives and design of the study.  If applicable, include a diagram or flowchart of the entire experimental design to illustrate the most important elements, such as specific materials, treatments, measurements, data collection, and methods of data analysis. This will facilitate the ability of editors, reviewers, and readers to understand and follow the concept of the study, the study design, and the results. 

\subsection{Subsection 1…n}
The rest of the Material and Methods should be divided by short subheadings for each method or technique. When a particularly lengthy description is required, a portion of the materials and methods can be included in the Supplementary Materials. This option should be used only in exceptional circumstances.

\subsection{Animal and Human Studies}
Studies involving animals or humans should include separate sections with the subheadings “Animals and Study Approval” or “Subjects and Study Approval,” as appropriate for animal or human research, respectively. All human studies must have been approved by the appropriate institutional review board(s). The Subjects and Study Approval subsection must include a specific declaration of such approval, including a statement indicating that written informed consent was received from participants prior to inclusion in the study. For animal models, the Animals and Study Approval subsection must include the precise genotype, strain, source, number of backcrosses, sex, and age of animals. Additionally, all animal studies must have been approved by the appropriate institutional review board(s). This subsection must include a specific declaration of such approval.

\subsection{Statistical Analysis}
The final subsection of the Materials and Methods should be titled “Statistical Analysis.”  This subsection describes the statistical methods with enough detail for verification of the results by a knowledgeable reader with access to the original data. Although this subsection describes the statistical methods used, the values for N, P, and the specific statistical test performed for each experiment should be included in the appropriate figure or table legend or main text. 

\section*{Acknowledgments}
Anyone who made a contribution to the research or manuscript, but who is not a listed author, should be acknowledged (with their permission). Types of acknowledgements include:

\subsection*{General} 
Thank others for any contributions, whether it be direct technical help or indirect assistance 

\subsection*{Author Contributions} 
Describe contributions of each author to the paper, using the first initial and full last name. 

\medskip Examples:

``S. Zhang conceived the idea and designed the experiments.''

``E. F. Mustermann and J. F. Smith conducted the experiments.''

``All authors contributed equally to the writing of the manuscript.''

\subsection*{Funding}
Name financially supporting bodies (written out in full), followed by the funding awardee and associated grant numbers (if applicable) in square brackets. 

\medskip Example: 

``This work was supported by the Engineering and Physical Sciences Research Council [grant numbers xxxx, yyyy]; the National Science Foundation [grant number zzzz]; and a Leverhulme Trust Research Project Grant.'' 

\medskip
If the research did not receive specific funding, but was performed as part of the employment of the authors, please name this employer. If the funder was involved in the manuscript writing, editing, approval, or decision to publish, please declare this.

\subsection*{Conflicts of Interest}
Conflicts of interest (COIs, also known as ``competing interests'') occur when issues outside research could be reasonably perceived to affect the neutrality or objectivity of the work or its assessment. 

Authors must declare all potential interests – whether or not they actually had an influence – in a ‘Conflicts of Interest’ section, which should explain why the interest may be a conflict. Authors must declare current or recent funding (including for Article Processing Charges) and other payments, goods or services that might influence the work. All funding, whether a conflict or not, must be declared in a ``Funding Statement.'' The involvement of anyone other than the authors who 1) has an interest in the outcome of the work; 2) is affiliated to an organization with such an interest; or 3) was employed or paid by a funder, in the commissioning, conception, planning, design, conduct, or analysis of the work, the preparation or editing of the manuscript, or the decision to publish must be declared.

If there are none, the authors should state ``The author(s) declare(s) that there is no conflict of interest regarding the publication of this article.'' Submitting authors are responsible for coauthors declaring their interests. Declared conflicts of interest will be considered by the editor and reviewers and included in the published article.

\subsection*{Data Availability}
A data availability statement is compulsory for all research articles. This statement describes whether and how others can access the data supporting the findings of the paper, including 1) what the nature of the data is, 2) where the data can be accessed, and 3) any restrictions on data access and why.

If data are in an archive, include the accession number or a placeholder for it. Also include any materials that must be obtained through a Material Transfer Agreements (MTA). 

\section*{Supplementary Materials}
Include supporting text, figures, and tables at the end of the main manuscript file if possible. Alternatively, Supplementary Materials can be included as a separate Word file. Include titles and captions for additional file types that cannot be embedded into the Word file. Supplementary figures should be embedded in the Word file in the order in which they are mentioned in the text, with the legends directly below the figure.
\begin{itemize}
    \item All should include a title in addition to a legend.
    \item Any references cited in the Supplementary Materials must already appear in the reference list; no separate supplementary reference list should be created.
    \item Supplementary Materials may include additional author notes — for example, a list of group authors.
    \item Supplementary Materials should be listed in the following order: supplementary text or materials and methods, supplementary figures, supplementary tables, other supplementary files (such as movies, data, interactive images, computer code, or database files), and references only cited in the supplementary materials. Be sure to submit all Supplementary Materials with the manuscript. Supplementary Materials should be named as follows:
    \end{itemize}
\medskip Example:

Supplementary Text

Fig. S1. Title of the first supplementary figure.

Fig. S2. Title of the second supplementary figure.

Table S1. Title of the first supplementary table.

Data file S1. Title of the first supplementary data file.

Movie S1. Title of the first supplementary movie.

Computer Code S1. Title of the first supplementary computer code.

Audio S1. Title of the first supplementary movie.

\section*{Guidelines for References}
There is only one reference list for all sources cited in the main text, figure and table legends, and Supplementary Materials. Do not include a second reference list in the Supplementary Materials section. References cited only in the Supplementary Materials section are not counted toward the limit of 40 references.
\begin{itemize}
    \item In the text, references should be numbered consecutively in the order of their first citation. Citations of references in the text should be identified using numbers in square brackets e.g., “as discussed by Liu [9]”; “as discussed elsewhere [9, 10]”. All references should be cited within the text and uncited references will be removed.
    \item Authors may submit the reference section in any numbered style for journal articles, as long as the style includes all authors (initials and last name), article title, journal title (or abbreviation), volume, year of publication, and pages. For journals that do not use page numbers, include the article number. For journals that do not use volume numbers, include the date of publication and DOI. If available, include a DOI for each reference.
    \item List all authors by first initial(s) and last name. Do not use op. cit., ibid., 3-m dashes, en dashes, or et al. (in place of the complete list of authors’ names).
    \item For online resources that do not have a DOI, include full URL, title of the page, source of the information, and date of access.
    \item If accepted, BME Frontiers will reformat the references in Chicago style. Authors are responsible for ensuring that the information in each reference is complete and accurate.
    \item Explanatory notes should be called out and cited separately and not be included as part of another cited reference. 
    \item Manuscripts should not include footnotes; information should be integrated into the text.  
\end{itemize}

\printbibliography

\end{document}
